\subsection{可微映射的限制}

设 n 与 r 是正整数，设 U 是 \(R^{n}\) 的开子集，并设 F 是巴拿赫空间。用 \(\mathcal{C}^{r}(\mathrm{U};\mathrm{F})\) 表示从 U 到 F 中的 \(C^{r}\) 类映射所成的向量空间，赋以 \(C^{r}\) 紧收敛拓扑。回忆它是当 K 取遍 U 的紧子集全体时，在 K 上的 \(C^{r}\) 一致收敛拓扑诸拓扑的上确界（VAR, R2, 12.3.10, p. 56）。空间 \(\mathcal{C}^{0}(\mathrm{U};\mathrm{F})\) 就是从 U 到 F 中的连续映射所成的空间 \(\mathcal{C}(\mathrm{U};\mathrm{F})\)，赋以紧收敛拓扑。

引理 3. --- 设 A 是由多重指标 \(\alpha \in \mathbf{N}^n\)（其中 \(|\alpha| \leqslant r\)）组成的集，并设 u 是线性映射 \(f \mapsto (\partial^\alpha f)_{\alpha \in \mathrm{A}}\)（它从 \(\mathcal{C}^r(\mathrm{U};\mathrm{F})\) 到 \(\mathcal{C}(\mathrm{U};\mathrm{F})^{\mathrm{A}}\) 中）。

\begin{enumerate}
\def\labelenumi{\alph{enumi})}
\item
  映射 u 是单射、连续的、严格的，且它的像是闭的。
\item
  拓扑向量空间 \(\mathcal{C}^{r}(\mathrm{U};\mathrm{F})\) 是完备的。
\end{enumerate}

映射 \(u\) 是线性的且是单射。由 \(\mathcal{C}^r (\mathrm{U};\mathrm{F})\) 的拓扑的定义，它是连续的且是严格的。

设 \((g_{\alpha})_{\alpha\in\mathrm{A}}\) 是 \(\mathcal{C}(\mathrm{U};\mathrm{F})^{\mathrm{A}}\) 中附着于 u 的像的一点。存在 \(\mathcal{C}^{r}(\mathrm{U};\mathrm{F})\) 上的一个滤子 F，使得 \(g_{\alpha}=\lim_{f,\mathcal{F}}\partial^{\alpha}f\) 在 \(\mathcal{C}(\mathrm{U};\mathrm{F})\) 中对一切 \(\alpha\in A\) 成立。设 m 是满足 \(0\leqslant m\leqslant r\) 的整数。对 m 用归纳法，由 FVR, II, p. 2 的 th. 1 可得，映射 \(g_{0}\) 是 \(C^{m}\) 类的，且 \(g_{\alpha}=\partial^{\alpha}g_{0}\) 对一切 \(\alpha\in\mathbf{N}^{n}\)（其中 \(|\alpha|\leqslant m\)）成立。于是 \(g_{0}\) 属于空间 \(\mathcal{C}^{r}(\mathrm{U};\mathrm{F})\)，且它在 u 下的像是 \((g_{\alpha})_{\alpha\in\mathrm{A}}\)。这就证明了 u 的像是闭的，从而建立了断言 a)。

断言 \(a\)）蕴含空间 \(\mathcal{C}^{r}(\mathrm{U};\mathrm{F})\) 同构于它在 \(\mathcal{C}(\mathrm{U};\mathrm{F})^{\mathrm{A}}\) 中的像；由于这个像是闭的，且空间 \(\mathcal{C}(\mathrm{U};\mathrm{F})^{\mathrm{A}}\) 是完备的（TG, X, p. 9, th. 2 的 cor. 3 以及 TG, II, p. 17, prop. 10），所以空间 \(\mathcal{C}^{r}(\mathrm{U};\mathrm{F})\) 是完备的。

命题 8. --- 假设 F 是有限维的。设 s 是满足 \(0 \leqslant s < r\) 的整数，并设 V 是 U 的相对紧开子集。线性映射 \(f \mapsto f|V\)（它从 \(\mathcal{C}^{r}(U;F)\) 到 \(\mathcal{C}^{s}(V;F)\) 中）是紧的。

设 W 是由这样的函数 \(f \in \mathcal{C}^{r}(\mathrm{U};\mathrm{F})\) 组成的集：其 \(\leqslant r\) 阶偏导数在 \(\overline{V}\) 的每一点处取范数小于 1 的值。集 W 是空间 \(\mathcal{C}^{r}(\mathrm{U};\mathrm{F})\) 中 0 的邻域。设 \(\alpha\) 是满足 \(|\alpha| \leqslant s\) 的多重指标。考虑由形如 \((\partial^{\alpha}f)|\mathrm{V}\)（其中 f 属于 W）的函数组成的集 H。由中值定理（VAR, R, 2.2.3），集 H 是 \(\mathcal{C}(\mathrm{V};\mathrm{F})\) 的等度连续子集。此外，对每个 \(x \in V\)，H 在映射 \(g \mapsto g(x)\) 下的像是 F 中的有界子集，因而是相对紧的。由阿斯科利定理（TG, X, p. 18, cor. 2），集 H 在 \(\mathcal{C}(\mathrm{V};\mathrm{F})\) 中是相对紧的。这就证明了线性映射 \(f \mapsto (\partial^{\alpha}f)|\mathrm{V}\)（它从 \(\mathcal{C}^{r}(\mathrm{U};\mathrm{F})\) 到 \(\mathcal{C}(\mathrm{V};\mathrm{F})\) 中）是紧的。命题于是由引理 3 以及 III, p. 3 的注记 5 和 6 得出。

